\chapter{选题依据}

\section{选题背景意义}
说明课题背景、研究必要性及其学术价值或应用价值。

\section{国内外研究现状}
围绕拟研究的问题梳理已有工作，比较研究所用的方法、证据与适用条件。

\subsection{研究方向一}
归纳具有代表性的研究。上标引文写法见示例\cite{sample-book}。

\MasterLevelFour{代表性工作}
四级标题使用括号编号。此处填写代表性工作的发现及其局限。

\subsection{研究方向二}
概括发展趋势，并说明尚待解决的问题如何引出本课题。

\chapter{研究目标与研究内容}

\section{研究目标}
列出拟达到的总体目标与具体目标。

\section{研究内容}
图~\ref{fig:proposal-roadmap}为技术路线示例。可修改图中文字和连接，或用 PDF、PNG 等替换。

\begin{figure}[htbp]
  \centering
  \begin{tikzpicture}[
    roadmap box/.style={draw, minimum width=27mm, minimum height=9mm, align=center},
    roadmap arrow/.style={-{Latex}, line width=0.5bp},
    node distance=8mm
  ]
    \node[roadmap box] (question) {界定问题};
    \node[roadmap box, right=of question] (method) {建立方法};
    \node[roadmap box, right=of method] (assessment) {验证与评价};
    \draw[roadmap arrow] (question) -- (method);
    \draw[roadmap arrow] (method) -- (assessment);
  \end{tikzpicture}
  \caption{技术路线示例}
  \label{fig:proposal-roadmap}
\end{figure}

\subsection{研究内容一}
说明具体任务、研究对象和拟交付的阶段结果。

\subsection{研究内容二}
说明各项任务的衔接关系及后续验证。

\section{拟解决的关键问题}
说明关键问题、当前难点和判断问题是否解决的依据。

\chapter{拟采用的研究方案}

\section{研究方法与实施步骤}
明确所采用的方法、数据或实验条件、关键实施步骤和对照设计。

\subsection{方法及参数定义}
例如，可用式~\eqref{eq:proposal-example}演示公式编号和物理量释义。
\begin{equation}
  y = f(x).
  \label{eq:proposal-example}
\end{equation}
\noindent 其中，$x$ 为输入量，$y$ 为输出量，$f$ 为映射函数。

\section{验证方案}
说明评价指标、对照条件及结果判定方法。

\chapter{可行性分析}

\section{已有基础与实施条件}
说明可以使用的数据、设备、人员与方法基础。

\section{关键风险与应对措施}
说明关键风险及相应的替代方案。

\chapter{预期目标与成果}

\section{预期目标}
写明开题后拟达到的定性或定量目标。

\section{预期成果}
列出预期形成的论文、软件、实验数据、原型或其他成果。

\chapter{论文进度安排}

论文进度安排见表~\ref{tab:proposal-schedule}，按实际执行时间填写。

\begin{table}[htbp]
  \centering
  \caption{学位论文时间进度安排表}
  \label{tab:proposal-schedule}
  \begin{MasterTable}
  \begin{tabularx}{0.9\textwidth}{@{}l>{\raggedright\arraybackslash}X@{}}
    \toprule
    时间 & 内容 \\
    \midrule
    YYYY年MM月--YYYY年MM月 & 文献调研、明确研究问题及评价方案。 \\
    YYYY年MM月--YYYY年MM月 & 实施研究方案，收集并整理实验结果。 \\
    YYYY年MM月--YYYY年MM月 & 完成分析、撰写及修改学位论文。 \\
    \bottomrule
  \end{tabularx}
  \end{MasterTable}
\end{table}
